\section{Latihan}
\label{sec:latihan-bab15}

\begin{enumerate}
    \item \textbf{Array satu dimensi.} Jelaskan cara mendeklarasikan
          array word lima elemen. Tentukan offset elemen pertama,
          kedua, dan kelima.

    \item \textbf{Array multi dimensi.} Diketahui matriks word
          $3 \times 4$. Tentukan offset elemen baris 2 kolom 1 dan
          baris 1 kolom 3.

    \item \textbf{String.} Jelaskan representasi string di memori.
          Mengapa string untuk layanan DOS fungsi \hex{09} harus
          diakhiri \code{'\$'}?

    \item \textbf{Menyalin string.} Tuliskan perulangan yang
          menyalin string dari \reg{SI} ke \reg{DI} sampai menemukan
          \code{'\$'}.

    \item \textbf{Membandingkan string.} Tuliskan perulangan yang
          membandingkan dua string karakter demi karakter. Kapan
          perbandingan dapat dihentikan?

    \item \textbf{Desimal.} Jelaskan algoritme mengubah bilangan
          menjadi digit desimal. Mengapa sisa pembagian disimpan
          pada \eng{stack}?

    \item \textbf{Heksadesimal.} Jelaskan algoritme mengubah
          bilangan menjadi empat digit heksadesimal. Mengapa nilai
          10--15 perlu ditambah 7 sebelum diubah menjadi karakter?

    \item \textbf{Penulisan program.} Tuliskan program yang
          menjumlahkan array word lima elemen, lalu mencetak hasilnya
          sebagai bilangan desimal. Gunakan \asm{LOOP} dan prosedur
          konversi desimal.
\end{enumerate}

\begin{refleksibox}
\vspace{1.1em}
Jawab pertanyaan berikut pada jurnal praktikum Anda.
\begin{enumerate}
    \item Bagian mana yang paling sulit: perhitungan offset array,
          penyalinan string, atau konversi bilangan? Mengapa?
    \item Setelah menelusuri konversi desimal, apakah Anda memahami
          mengapa \eng{stack} dipakai untuk membalik urutan digit?
          Jelaskan.
    \item Konsep apa yang perlu Anda ulangi sebelum mempelajari
          debugging level instruksi pada bab berikutnya?
\end{enumerate}
\end{refleksibox}
